% AI-produced by GPT-6 Astra (OpenAI, ChatGPT Pro mode).
\subsection{The comparison fibre in base dimension four}

\begin{saresult}{A.135}{Proved under stated hypotheses}{An exact comparison-fibre reduction for closed four-dimensional bases}
Let $B$ be a closed compact real PL four-manifold. All PL families below
are jointly PL and preserve the parameter projection; all TOP spaces
are replaced by their singular complexes. Set
\[
\begin{aligned}
 A_C&=\operatorname{Aut}_C(B),\\
 P_C&=\operatorname{Aut}_C(B\times I\text{ rel }B\times\{0\}),\\
 H_C&=\operatorname{Aut}_C(B\times I\text{ rel }B\times\partial I),
 \qquad C\in\{\mathrm{PL},\mathrm{TOP}\},
\end{aligned}
\]
and let $F_A,F_P,F_H$ be the based homotopy fibres of the respective
natural maps from PL to TOP, taken over the identity.
Top restriction induces a map $F_P\to F_A$ with homotopy fibre
\[
 F_H\simeq \operatorname{Map}(B,K(\mathbb Z/2,1)).
                                                        \tag{A.135.1}
\]
Here the mapping space is unbased, pointed by the constant map.
Consequently there is a based homotopy-fibre sequence
\[
 \operatorname{Map}(B,K(\mathbb Z/2,1))
 \longrightarrow F_P\longrightarrow F_A.                \tag{A.135.2}
\]
In particular $\pi_jF_P\to\pi_jF_A$ is an isomorphism for $j\ge3$,
and there is the exact segment
\[
\begin{aligned}
0&\longrightarrow\pi_2F_P\longrightarrow\pi_2F_A
\longrightarrow H^0(B;\mathbb Z/2)\\
 &\longrightarrow\pi_1F_P\longrightarrow\pi_1F_A
\longrightarrow H^1(B;\mathbb Z/2)\\
 &\longrightarrow\pi_0F_P\longrightarrow\pi_0F_A.
\end{aligned}                                           \tag{A.135.3}
\]
The last line is understood with its usual pointed-set exactness.
This identifies the fully boundary-fixed comparison fibre and reduces
the free-top comparison to the ordinary four-manifold automorphism
comparison. It does not assert that $F_A$ or $F_P$ is contractible,
that the original comparison is surjective on components, or that
all relative PL structures have the original PL isomorphism type.
The connecting maps and any resulting extension data in (A.135.3)
remain unevaluated; the abstract homotopy type of $F_A$ alone does
not determine $F_P$.
\end{saresult}
\begin{proof}
If $B$ is empty, every space in the statement is a point, so suppose
that $B$ is nonempty.
\textbf{Top restriction is the required ordinary-family fibration.}
For either category the sequence
\[
 H_C\longrightarrow P_C\xrightarrow{e} A_C             \tag{A.135.4}
\]
is a based homotopy-fibre sequence. We verify the precise lifting
property instead of using a block concordance space.
Let $J=\Lambda_i^q\subset\Delta^q$ be a horn, let $p_J$ be a
$J$-family in $P_C$, and let $a$ be a $\Delta^q$-family in $A_C$
whose restriction to $J$ is $e(p_J)$. Choose a PL retraction
$r:\Delta^q\to J$. There is a jointly PL homotopy
$d:\Delta^q\times I\to\Delta^q$ with
$d_0=r$, $d_1=\operatorname{id}$ and $d_s|J=\operatorname{id}$.
Indeed triangulate the prism compatibly with $r$ on its bottom and
with $J\times I$; prescribe $r$ on the bottom, the identity on the
top and the projection on $J\times I$, then extend linearly on a
refinement. Convexity of $\Delta^q$ keeps the extension in the target.
No homeomorphism condition is needed on $d_s$.

Put $p^0_t=p_{r(t)}$ and
$\delta_t=a_t\,e(p^0_t)^{-1}$. Thus $\delta_t=1$ on $J$.
On $B\times I$ define
\[
 E_t(x,u)=(\delta_{d_u(t)}(x),u),\qquad
 p_t=E_t p^0_t.                                         \tag{A.135.5}
\]
The bottom restriction of $E_t$ is the identity because $d_0(t)\in J$;
its top restriction is $\delta_t$; and it is the identity for $t\in J$.
It follows that $p$ fills the horn and $e(p)=a$ exactly.
The total map is jointly PL when the input is: composition with the
PL map $(t,u)\mapsto d_u(t)$ preserves that property, and its inverse
uses the corresponding total inverse family of $\delta$. The same
formula is continuous in TOP. Thus $e$ is a Kan fibration, with literal
kernel $H_C$. No surjectivity on the set of vertices or on all
components of $A_C$ is claimed or needed.

Apply the homotopy-fibre construction vertically to the natural
comparison of (A.135.4) in the two categories. Homotopy fibres commute
with homotopy fibres: both iterated constructions are the homotopy
limit of the same square with its specified identity basepoints.
Since each row has homotopy fibre $H_C$, this gives
\[
 F_H\longrightarrow F_P\longrightarrow F_A.             \tag{A.135.6}
\]
The statement is a homotopy-fibre sequence; it does not require the
map $F_P\to F_A$ in a particular path model to be a strict fibration
surjective on vertices.

\textbf{Pointwise boundary fixing versus fixing a collar.}
Write $N=B\times I$. We will use the relative classification theorem
with a neighbourhood fixed. Here is an explicit comparison with the
pointwise-fixed groups in the statement. Choose $0<\delta<1/2$.
For $0\le s\le\delta/2$ put
\[
 \sigma_s(t)=\min\{t+s,(t+\delta)/2\},\qquad 0\le t\le\delta.
                                                        \tag{A.135.7}
\]
This is a strictly increasing PL embedding $[0,\delta]\to[s,\delta]$,
equal to the identity at $s=0$ and fixing $\delta$. Its total map and
its inverse on the total image are PL. Use this formula in the bottom
collar and the reflected formula in the top collar, with the identity
between them, to obtain embeddings $c_s:N\to N$ with image
$B\times[s,1-s]$.
For $h\in H_C$ define
\[
 R_s(h)=
 \begin{cases}
 c_s h c_s^{-1}&\text{on }c_s(N),\\
 \operatorname{id}&\text{on }N\setminus\operatorname{int}c_s(N).
 \end{cases}                                           \tag{A.135.8}
\]
The two definitions agree on the gluing boundary, since $h$ fixes
both boundary faces pointwise. The formula gives continuous TOP
families and jointly PL families in the PL case. It starts at $h$;
for every fixed $s>0$ it fixes a uniform collar.
Moreover $\sigma_s(t)\ge t$. Hence a family already fixed on a
collar stays fixed on that collar throughout this deformation.
The identity family stays the identity, including on any prescribed
parameter subcomplex. Applying (A.135.8) to finite representatives of
homotopy groups and their relative nullhomotopies proves that the
inclusion of the collar-fixed complex $H_C^{\mathrm{col}}\to H_C$
is a weak equivalence in both categories. Here a simplex of
$H_C^{\mathrm{col}}$ fixes one common collar for its whole parameter
simplex; the collar may shrink when passing to a different simplex.
A finite family of simplices admits a common smaller collar. This is
the direct-limit complex with this convention, not an unproved
identification with the singular complex of a union of topological
groups. This argument also provides a common collar for every finite
relative nullhomotopy. The comparison
square commutes, so it preserves the homotopy fibre $F_H$.

\textbf{The relative structure space, with all its boundary controlled.}
Attach infinite exterior collars to $N$ and put
\[
 M=B\times\mathbb R,\qquad
 C=M\setminus(B\times(0,1)).
\]
Then $M$ is an open PL five-manifold with no boundary and $C$ is closed.
Extension by the identity identifies $H_C^{\mathrm{col}}$ (where here
the subscript denotes the category) with the complex of automorphisms
of $M$ fixed on some neighbourhood of this closed set $C$.
The compact support in $N$ ensures that the locally finite PL convention
on $M$ restricts to exactly the finite joint PL convention on $N$.
Let $\mathcal S$ be the simplicial space of PL structures on $M$,
sliced over the parameter simplex, equal to the standard structure
$\alpha$ near $C$.

Kirby--Siebenmann, Essay~V, (1.5) and its relative version on p.~229,
give the quotient map
\[
 \operatorname{Aut}_{\mathrm{TOP}}(M\text{ near }C)
 \longrightarrow
 \operatorname{Aut}_{\mathrm{TOP}}(M\text{ near }C)
 /\operatorname{Aut}_{\mathrm{PL}}(M\text{ near }C)
                                                        \tag{A.135.9}
\]
as a Kan fibration. Its target embeds in $\mathcal S$ with image a union
of components containing the component $\mathcal S_\alpha$ of the
standard structure \cite{KirbySiebenmann1977}. These are actual
projection-preserving families. The hypotheses hold because the
ambient manifold is the boundaryless five-manifold $M$. The relative
version fixes a neighbourhood of $C$, precisely as arranged above.
In particular, taking based loops loses none of the relevant structure
space:
\[
 F_H\simeq\Omega\mathcal S_\alpha.                       \tag{A.135.10}
\]
Indeed the homotopy fibre of the inclusion of the fibre in the total
space of (A.135.9) is the loop space of its base at the identity coset.
The base component is $\mathcal S_\alpha$. This is the only orbit
assertion used; other components of $\mathcal S$ need not be in this
orbit.

The relative classification (3.2), Essay~V, pp.~240--241, identifies
$\mathcal S$ with the space of lifts of the tangent classifying map
$M\to B\mathrm{TOP}(5)$ to $B\mathrm{PL}(5)$, fixed near $C$ by the
standard PL lift. The relative restriction map is a fibration for
$m\ge5$, without an additional cohomological condition. Theorem~5.5(I),
p.~251, identifies the fibre $\mathrm{TOP}(5)/\mathrm{PL}(5)$ with
$K(\mathbb Z/2,3)$ and identifies stabilization with a weak equivalence
on these fibres.

There is no unproved untwisting assertion here. Essay~V, \S5.6,
pp.~252--253, realizes the stable lifting problem using the universal
Kirby--Siebenmann map to $K(\mathbb Z/2,4)$. The standard PL lift gives
a chosen nullhomotopy of its pullback. Subtracting this nullhomotopy
identifies the lift space, relative to the same lift near $C$, with
maps to the loop space $K(\mathbb Z/2,3)$ that vanish near $C$.
The relative form is also obtained directly by the natural square
(3.1)--(3.2). Thus
\[
 \mathcal S\simeq
 \operatorname{Map}_*(N/\partial N,K(\mathbb Z/2,3)).      \tag{A.135.11}
\]
Replacing maps constant on a neighbourhood of the boundary by maps
constant on the boundary is harmless: collar compression, fixed on the
boundary and relative to any already constant parameter faces, gives
a deformation to a uniformly constant collar for each finite family.
Equivalently this is the standard mapping-space comparison for the
collared cofibration $\partial N\subset N$.
The chosen standard lift corresponds to the constant map.

Since $N/\partial N\cong\Sigma B_+$, where $B_+$ means $B$ with a
disjoint basepoint, adjunction gives
\[
 \begin{aligned}
 \Omega\mathcal S_\alpha
 &\simeq\Omega\operatorname{Map}_*(\Sigma B_+,K(\mathbb Z/2,3))\\
 &\simeq\operatorname{Map}_*(B_+,\Omega^2K(\mathbb Z/2,3))\\
 &\simeq\operatorname{Map}(B,K(\mathbb Z/2,1)).
 \end{aligned}                                         \tag{A.135.12}
\]
The first loop ignores the other components, as it should. This proves
(A.135.1), and (A.135.6) proves (A.135.2).
For this last mapping space the components are $H^1(B;\mathbb Z/2)$,
its fundamental group at the constant map is $H^0(B;\mathbb Z/2)$,
and its higher homotopy groups vanish. The long exact sequence gives
(A.135.3) and the asserted isomorphisms in degrees at least three.

\textbf{Why this does not remove the dimensional gate.}
The free top has survived as $F_A$ in (A.135.2). None of the argument
classifies the PL structures on the compact four-manifold $B$ itself,
or proves surjectivity on the components of the original free-top
comparison. Nor can the same proof simply be used for a closed
three-manifold base. Then the open ambient manifold has dimension
four, and the relative statement on p.~240 requires
$H^4(M,C;\mathbb Z/2)=0$. Here this group is
$H^3(B;\mathbb Z/2)$, which is nonzero for every nonempty closed
connected three-manifold. The hypothesis therefore fails in exactly
that proposed application.
\end{proof}
